
\documentclass{article}
\usepackage{colortbl}
\usepackage{makecell}
\usepackage{multirow}
\usepackage{supertabular}

\begin{document}

\newcounter{utterance}

\centering \large Interaction Transcript for game `clean\_up', experiment `1\_medium\_3obj\_sv', episode 0 with deepseek{-}v4{-}pro.
\vspace{24pt}

{ \footnotesize  \setcounter{utterance}{1}
\setlength{\tabcolsep}{0pt}
\begin{supertabular}{c@{$\;$}|p{.15\linewidth}@{}p{.15\linewidth}p{.15\linewidth}p{.15\linewidth}p{.15\linewidth}p{.15\linewidth}}
    \# & $\;$A & \multicolumn{4}{c}{Game Master} & $\;\:$B\\
    \hline

    \theutterance \stepcounter{utterance}  
    & & \multicolumn{4}{p{0.6\linewidth}}{
        \cellcolor[rgb]{0.9,0.9,0.9}{
            \makecell[{{p{\linewidth}}}]{
                \texttt{\tiny{[P1$\langle$GM]}}
                \texttt{Jag är din spelledare och du spelar ett samarbetsspel med följande rutnät som spelplan:} \\
\\ 
\texttt{```} \\
\texttt{ 1234567} \\
\texttt{╔════╤══╗} \\
\texttt{║◌C◌◌│◌◌║ 1} \\
\texttt{║◌◌L◌├──╢ 2} \\
\texttt{║◌◌◌P│◌◌║ 3} \\
\texttt{║◌┌──┼──╢ 4} \\
\texttt{║◌│◌◌│◌◌║ 5} \\
\texttt{╟─┴┬─┤◌◌║ 6} \\
\texttt{║◌◌│◌│◌◌║ 7} \\
\texttt{╚══╧═╧══╝} \\
\\ 
\texttt{```} \\
\\ 
\texttt{* Den övre kanten visar x{-}koordinater som ökar åt höger, och den högra kanten y{-}koordinater som ökar nedåt.} \\
\texttt{* Följande objekt är slumpmässigt placerade på ditt rutnät: 'C', 'L', 'P'.} \\
\\ 
\\ 
\texttt{Den andra spelaren ser en variant av spelplanen, där objekten är placerade på olika slumpmässiga platser. Du kan inte se den andra spelarens plan, och de kan inte se din.} \\
\\ 
\\ 
\texttt{**Mål:**} \\
\\ 
\texttt{Båda spelarna måste flytta objekten på sin respektive bakgrund så att identiska objekt hamnar på samma koordinater. Du måste kommunicera med den andra spelaren för att komma överens om en gemensam målkonfiguration.} \\
\\ 
\\ 
\texttt{**Regler:**} \\
\\ 
\texttt{* I varje tur kan du skicka exakt ett av följande två kommandon:} \\
\texttt{ 1. `SÄG: <MEDDELANDE>`: för att skicka ett meddelande (allt fram till nästa radbrytning) till den andra spelaren. Jag kommer att vidarebefordra det till din partner.} \\
\texttt{2. `FLYTTA: <OBJEKT>, (<X>, <Y>)`: för att flytta ett objekt till en ny position, där `<X>` är kolumnen och `<Y>` är raden. Jag kommer att informera dig om ditt drag var framgångsrikt eller inte.} \\
\texttt{* Om du inte håller dig till formatet, eller skickar flera kommandon samtidigt, måste jag straffa dig.} \\
\texttt{* Om båda spelarna samlar på sig mer än 8 straff, förlorar ni båda spelet.} \\
\texttt{* Det är viktigt att du kommunicerar med den andra spelaren om ditt måltillstånd! Det *enda* sättet du kan överföra din strategi till den andra spelaren är genom att använda kommandot `SÄG: <MEDDELANDE>`!} \\
\\ 
\\ 
\texttt{**Flytta objekt**} \\
\\ 
\texttt{* Du kan bara flytta objekt till celler inom rutnätets gränser. Målcellen måste vara tom, d.v.s. den får bara innehålla symbolen '◌'.} \\
\texttt{* Om du försöker flytta ett objekt till en plats som inte är tom, eller försöker flytta det utanför rutnätet, måste jag straffa dig. Du får göra ett nytt försök.} \\
\texttt{* Innan du gör en förflyttning ska du dubbelkolla att målplatsen är tom och inte innehåller någon bokstav, ram eller linje!} \\
\\ 
\\ 
\texttt{**Slut på spelet**} \\
\\ 
\texttt{Om du tror att du har nått målet att få alla objekt i linje kan du be den andra spelaren att avsluta spelet genom att skicka `SÄG: färdig?`. Om den andra spelaren ber dig att avsluta spelet och du svarar med `SÄG: färdig!`, avslutas spelet.} \\
\\ 
\texttt{Båda spelarna vinner om spelet avslutas inom 12 rundor, där en runda definieras som att två spelare var skickar ett giltigt kommando.} \\
\\ 
\texttt{**Poäng:**} \\
\\ 
\texttt{Ju närmare de identiska objekten är på båda spelplanerna, desto fler poäng får du. Straff minskar dina poäng. Kan du slå rekordet?} \\
\\ 
\\ 
\texttt{Skicka ett meddelande till den andra spelaren för att starta spelet!} \\
            }
        }
    }
    & \\ \\

    \theutterance \stepcounter{utterance}  
    & \multicolumn{4}{p{0.6\linewidth}}{
        \cellcolor[rgb]{0.9,0.9,0.9}{
            \makecell[{{p{\linewidth}}}]{
                \texttt{\tiny{[P1$\rangle$GM]}}
                \texttt{SÄG: Hej! Jag har C på (2,1), L på (3,2) och P på (4,3) på min spelplan. Var har du dina objekt?} \\
            }
        }
    }
    & & \\ \\

    \theutterance \stepcounter{utterance}  
    & & & \multicolumn{2}{p{0.3\linewidth}}{
        \cellcolor[rgb]{0.9,0.9,0.9}{
            \makecell[{{p{\linewidth}}}]{
                \texttt{\tiny{[GM$|$GM]}}
                \texttt{SÄG: Hej! Jag har C på (2,1), L på (3,2) och P på (4,3) på min spelplan. Var har du dina objekt?} \\
            }
        }
    }
    & & \\ \\

    \theutterance \stepcounter{utterance}  
    & & \multicolumn{4}{p{0.6\linewidth}}{
        \cellcolor[rgb]{0.9,0.9,0.9}{
            \makecell[{{p{\linewidth}}}]{
                \texttt{\tiny{[GM$\rangle$P2]}}
                \texttt{Jag är din spelledare och du spelar ett samarbetsspel med följande rutnät som spelplan:} \\
\\ 
\texttt{```} \\
\texttt{ 1234567} \\
\texttt{╔════╤══╗} \\
\texttt{║◌◌P◌│◌◌║ 1} \\
\texttt{║◌◌C◌├──╢ 2} \\
\texttt{║◌◌◌◌│◌◌║ 3} \\
\texttt{║◌┌──┼──╢ 4} \\
\texttt{║◌│◌◌│◌◌║ 5} \\
\texttt{╟─┴┬─┤◌◌║ 6} \\
\texttt{║L◌│◌│◌◌║ 7} \\
\texttt{╚══╧═╧══╝} \\
\\ 
\texttt{```} \\
\\ 
\texttt{* Den övre kanten visar x{-}koordinater som ökar åt höger, och den högra kanten y{-}koordinater som ökar nedåt.} \\
\texttt{* Följande objekt är slumpmässigt placerade på ditt rutnät: 'C', 'L', 'P'.} \\
\\ 
\\ 
\texttt{Den andra spelaren ser en variant av spelplanen, där objekten är placerade på olika slumpmässiga platser. Du kan inte se den andra spelarens plan, och de kan inte se din.} \\
\\ 
\\ 
\texttt{**Mål:**} \\
\\ 
\texttt{Båda spelarna måste flytta objekten på sin respektive bakgrund så att identiska objekt hamnar på samma koordinater. Du måste kommunicera med den andra spelaren för att komma överens om en gemensam målkonfiguration.} \\
\\ 
\\ 
\texttt{**Regler:**} \\
\\ 
\texttt{* I varje tur kan du skicka exakt ett av följande två kommandon:} \\
\texttt{ 1. `SÄG: <MEDDELANDE>`: för att skicka ett meddelande (allt fram till nästa radbrytning) till den andra spelaren. Jag kommer att vidarebefordra det till din partner.} \\
\texttt{2. `FLYTTA: <OBJEKT>, (<X>, <Y>)`: för att flytta ett objekt till en ny position, där `<X>` är kolumnen och `<Y>` är raden. Jag kommer att informera dig om ditt drag var framgångsrikt eller inte.} \\
\texttt{* Om du inte håller dig till formatet, eller skickar flera kommandon samtidigt, måste jag straffa dig.} \\
\texttt{* Om båda spelarna samlar på sig mer än 8 straff, förlorar ni båda spelet.} \\
\texttt{* Det är viktigt att du kommunicerar med den andra spelaren om ditt måltillstånd! Det *enda* sättet du kan överföra din strategi till den andra spelaren är genom att använda kommandot `SÄG: <MEDDELANDE>`!} \\
\\ 
\\ 
\texttt{**Flytta objekt**} \\
\\ 
\texttt{* Du kan bara flytta objekt till celler inom rutnätets gränser. Målcellen måste vara tom, d.v.s. den får bara innehålla symbolen '◌'.} \\
\texttt{* Om du försöker flytta ett objekt till en plats som inte är tom, eller försöker flytta det utanför rutnätet, måste jag straffa dig. Du får göra ett nytt försök.} \\
\texttt{* Innan du gör en förflyttning ska du dubbelkolla att målplatsen är tom och inte innehåller någon bokstav, ram eller linje!} \\
\\ 
\\ 
\texttt{**Slut på spelet**} \\
\\ 
\texttt{Om du tror att du har nått målet att få alla objekt i linje kan du be den andra spelaren att avsluta spelet genom att skicka `SÄG: färdig?`. Om den andra spelaren ber dig att avsluta spelet och du svarar med `SÄG: färdig!`, avslutas spelet.} \\
\\ 
\texttt{Båda spelarna vinner om spelet avslutas inom 12 rundor, där en runda definieras som att två spelare var skickar ett giltigt kommando.} \\
\\ 
\texttt{**Poäng:**} \\
\\ 
\texttt{Ju närmare de identiska objekten är på båda spelplanerna, desto fler poäng får du. Straff minskar dina poäng. Kan du slå rekordet?} \\
\\ 
\texttt{Den andra spelaren startade spelet genom att skicka detta meddelande:} \\
\\ 
\texttt{"Hej! Jag har C på (2,1), L på (3,2) och P på (4,3) på min spelplan. Var har du dina objekt?"} \\
\\ 
\texttt{Vad är ditt första kommando?} \\
            }
        }
    }
    & \\ \\

    \theutterance \stepcounter{utterance}  
    & & & \multicolumn{4}{p{0.6\linewidth}}{
        \cellcolor[rgb]{0.9,0.9,0.9}{
            \makecell[{{p{\linewidth}}}]{
                \texttt{\tiny{[GM$\langle$P2]}}
                \texttt{Jag börjar med att undersöka min spelplan för att hitta var mina objekt finns.} \\
\\ 
\texttt{Min spelplan:} \\
\\ 
\texttt{{-} C finns på (3,2)} \\
\texttt{{-} L finns på (1,7)} \\
\texttt{{-} P finns på (3,1)} \\
\\ 
\texttt{Jämfört med den andra spelarens positioner:} \\
\texttt{{-} De har C på (2,1), jag har C på (3,2)} \\
\texttt{{-} De har L på (3,2), jag har L på (1,7)} \\
\texttt{{-} De har P på (4,3), jag har P på (3,1)} \\
\\ 
\texttt{Jag ska svara med att berätta var mina objekt finns.} \\
\\ 
\texttt{SÄG: Hej! Jag har C på (3,2), L på (1,7) och P på (3,1) på min spelplan. Vi behöver komma överens om gemensamma målpositioner. Har du något förslag?} \\
            }
        }
    }
     \\ \\

    \theutterance \stepcounter{utterance}  
    & & & \multicolumn{2}{p{0.3\linewidth}}{
        \cellcolor[rgb]{0.9,0.9,0.9}{
            \makecell[{{p{\linewidth}}}]{
                \texttt{\tiny{[GM$|$GM]}}
                \texttt{Jag börjar med att undersöka min spelplan för att hitta var mina objekt finns.} \\
\\ 
\texttt{Min spelplan:} \\
\\ 
\texttt{{-} C finns på (3,2)} \\
\texttt{{-} L finns på (1,7)} \\
\texttt{{-} P finns på (3,1)} \\
\\ 
\texttt{Jämfört med den andra spelarens positioner:} \\
\texttt{{-} De har C på (2,1), jag har C på (3,2)} \\
\texttt{{-} De har L på (3,2), jag har L på (1,7)} \\
\texttt{{-} De har P på (4,3), jag har P på (3,1)} \\
\\ 
\texttt{Jag ska svara med att berätta var mina objekt finns.} \\
\\ 
\texttt{SÄG: Hej! Jag har C på (3,2), L på (1,7) och P på (3,1) på min spelplan. Vi behöver komma överens om gemensamma målpositioner. Har du något förslag?} \\
            }
        }
    }
    & & \\ \\

    \theutterance \stepcounter{utterance}  
    & & & \multicolumn{2}{p{0.3\linewidth}}{
        \cellcolor[rgb]{0.9,0.9,0.9}{
            \makecell[{{p{\linewidth}}}]{
                \texttt{\tiny{[GM$|$GM]}}
                \texttt{Invalid format: head is not empty: 'Jag börjar med att undersöka min spelplan för att hitta var mina objekt finns.} \\
\\ 
\texttt{Min spelplan:} \\
\\ 
\texttt{{-} C finns på (3,2)} \\
\texttt{{-} L finns på (1,7)} \\
\texttt{{-} P finns på (3,1)} \\
\\ 
\texttt{Jämfört med den andra spelarens positioner:} \\
\texttt{{-} De har C på (2,1), jag har C på (3,2)} \\
\texttt{{-} De har L på (3,2), jag har L på (1,7)} \\
\texttt{{-} De har P på (4,3), jag har P på (3,1)} \\
\\ 
\texttt{Jag ska svara med att berätta var mina objekt finns.} \\
\\ 
\texttt{'} \\
            }
        }
    }
    & & \\ \\

    \theutterance \stepcounter{utterance}  
    & & \multicolumn{4}{p{0.6\linewidth}}{
        \cellcolor[rgb]{0.9,0.9,0.9}{
            \makecell[{{p{\linewidth}}}]{
                \texttt{\tiny{[GM$\rangle$P2]}}
                \texttt{Straff: Ditt meddelande får inte innehålla något före kommandot!} \\
\\ 
\texttt{Se till att ditt svar endast innehåller antingen `SÄG: <MEDDELANDE>` eller `FLYTTA: <OBJEKT>, (<X>, <Y>)`, och inget annat!} \\
\\ 
\texttt{Ni har tillsammans ackumulerat 1 av 8 straff.} \\
\texttt{Vänligen försök igen!} \\
            }
        }
    }
    & \\ \\

    \theutterance \stepcounter{utterance}  
    & & & \multicolumn{4}{p{0.6\linewidth}}{
        \cellcolor[rgb]{0.9,0.9,0.9}{
            \makecell[{{p{\linewidth}}}]{
                \texttt{\tiny{[GM$\langle$P2]}}
                \texttt{SÄG: Hej! Jag har C på (3,2), L på (1,7) och P på (3,1) på min spelplan. Vi behöver komma överens om gemensamma målpositioner. Har du något förslag?} \\
            }
        }
    }
     \\ \\

    \theutterance \stepcounter{utterance}  
    & & & \multicolumn{2}{p{0.3\linewidth}}{
        \cellcolor[rgb]{0.9,0.9,0.9}{
            \makecell[{{p{\linewidth}}}]{
                \texttt{\tiny{[GM$|$GM]}}
                \texttt{SÄG: Hej! Jag har C på (3,2), L på (1,7) och P på (3,1) på min spelplan. Vi behöver komma överens om gemensamma målpositioner. Har du något förslag?} \\
            }
        }
    }
    & & \\ \\

    \theutterance \stepcounter{utterance}  
    & & \multicolumn{4}{p{0.6\linewidth}}{
        \cellcolor[rgb]{0.9,0.9,0.9}{
            \makecell[{{p{\linewidth}}}]{
                \texttt{\tiny{[P1$\langle$GM]}}
                \texttt{Ditt meddelande har vidarebefordrats till den andra spelaren.} \\
\\ 
\\ 
\texttt{Du spelar för närvarande runda 1 av maximalt 12.} \\
\texttt{Ni har tillsammans ackumulerat 1 av 8 straff.} \\
\texttt{Den andra spelaren skickade detta meddelande:} \\
\\ 
\\ 
\texttt{"Hej! Jag har C på (3,2), L på (1,7) och P på (3,1) på min spelplan. Vi behöver komma överens om gemensamma målpositioner. Har du något förslag?"} \\
\\ 
\\ 
\texttt{Vad är ditt nästa kommando?} \\
            }
        }
    }
    & \\ \\

    \theutterance \stepcounter{utterance}  
    & \multicolumn{4}{p{0.6\linewidth}}{
        \cellcolor[rgb]{0.9,0.9,0.9}{
            \makecell[{{p{\linewidth}}}]{
                \texttt{\tiny{[P1$\rangle$GM]}}
                \texttt{SÄG: Okej, låt oss försöka hitta en gemensam lösning. Ett förslag: vi kan flytta C till (3,1), L till (1,7) och P till (4,3). Då matchar C din position, L din position, och P min position. Vad tror du om det? Eller har du ett bättre förslag?} \\
            }
        }
    }
    & & \\ \\

    \theutterance \stepcounter{utterance}  
    & & & \multicolumn{2}{p{0.3\linewidth}}{
        \cellcolor[rgb]{0.9,0.9,0.9}{
            \makecell[{{p{\linewidth}}}]{
                \texttt{\tiny{[GM$|$GM]}}
                \texttt{SÄG: Okej, låt oss försöka hitta en gemensam lösning. Ett förslag: vi kan flytta C till (3,1), L till (1,7) och P till (4,3). Då matchar C din position, L din position, och P min position. Vad tror du om det? Eller har du ett bättre förslag?} \\
            }
        }
    }
    & & \\ \\

    \theutterance \stepcounter{utterance}  
    & & \multicolumn{4}{p{0.6\linewidth}}{
        \cellcolor[rgb]{0.9,0.9,0.9}{
            \makecell[{{p{\linewidth}}}]{
                \texttt{\tiny{[GM$\rangle$P2]}}
                \texttt{Ditt meddelande har vidarebefordrats till den andra spelaren.} \\
\\ 
\\ 
\texttt{Du spelar för närvarande runda 2 av maximalt 12.} \\
\texttt{Ni har tillsammans ackumulerat 1 av 8 straff.} \\
\texttt{Den andra spelaren skickade detta meddelande:} \\
\\ 
\\ 
\texttt{"Okej, låt oss försöka hitta en gemensam lösning. Ett förslag: vi kan flytta C till (3,1), L till (1,7) och P till (4,3). Då matchar C din position, L din position, och P min position. Vad tror du om det? Eller har du ett bättre förslag?"} \\
\\ 
\\ 
\texttt{Vad är ditt nästa kommando?} \\
            }
        }
    }
    & \\ \\

    \theutterance \stepcounter{utterance}  
    & & & \multicolumn{4}{p{0.6\linewidth}}{
        \cellcolor[rgb]{0.9,0.9,0.9}{
            \makecell[{{p{\linewidth}}}]{
                \texttt{\tiny{[GM$\langle$P2]}}
                \texttt{SÄG: Bra förslag! Men jag kan inte flytta C till (3,1) eftersom P redan står där på min plan. Vi kan istället sikta på: C till (3,3), L till (2,4), P till (4,3). Då får vi ett gemensamt ledigt område i mitten. Vad tror du om det?} \\
            }
        }
    }
     \\ \\

    \theutterance \stepcounter{utterance}  
    & & & \multicolumn{2}{p{0.3\linewidth}}{
        \cellcolor[rgb]{0.9,0.9,0.9}{
            \makecell[{{p{\linewidth}}}]{
                \texttt{\tiny{[GM$|$GM]}}
                \texttt{SÄG: Bra förslag! Men jag kan inte flytta C till (3,1) eftersom P redan står där på min plan. Vi kan istället sikta på: C till (3,3), L till (2,4), P till (4,3). Då får vi ett gemensamt ledigt område i mitten. Vad tror du om det?} \\
            }
        }
    }
    & & \\ \\

    \theutterance \stepcounter{utterance}  
    & & \multicolumn{4}{p{0.6\linewidth}}{
        \cellcolor[rgb]{0.9,0.9,0.9}{
            \makecell[{{p{\linewidth}}}]{
                \texttt{\tiny{[P1$\langle$GM]}}
                \texttt{Ditt meddelande har vidarebefordrats till den andra spelaren.} \\
\\ 
\\ 
\texttt{Du spelar för närvarande runda 2 av maximalt 12.} \\
\texttt{Ni har tillsammans ackumulerat 1 av 8 straff.} \\
\texttt{Den andra spelaren skickade detta meddelande:} \\
\\ 
\\ 
\texttt{"Bra förslag! Men jag kan inte flytta C till (3,1) eftersom P redan står där på min plan. Vi kan istället sikta på: C till (3,3), L till (2,4), P till (4,3). Då får vi ett gemensamt ledigt område i mitten. Vad tror du om det?"} \\
\\ 
\\ 
\texttt{Vad är ditt nästa kommando?} \\
            }
        }
    }
    & \\ \\

    \theutterance \stepcounter{utterance}  
    & \multicolumn{4}{p{0.6\linewidth}}{
        \cellcolor[rgb]{0.9,0.9,0.9}{
            \makecell[{{p{\linewidth}}}]{
                \texttt{\tiny{[P1$\rangle$GM]}}
                \texttt{SÄG: Låt mig kolla min spelplan. Jag har tomma rutor på (3,3), (2,4) och (4,3) så det fungerar för mig! Jag kör på ditt förslag: C till (3,3), L till (2,4), P till (4,3). Jag börjar flytta mina objekt nu. Ska vi säga att vi meddelar varandra när vi är klara med alla flyttar?} \\
            }
        }
    }
    & & \\ \\

    \theutterance \stepcounter{utterance}  
    & & & \multicolumn{2}{p{0.3\linewidth}}{
        \cellcolor[rgb]{0.9,0.9,0.9}{
            \makecell[{{p{\linewidth}}}]{
                \texttt{\tiny{[GM$|$GM]}}
                \texttt{SÄG: Låt mig kolla min spelplan. Jag har tomma rutor på (3,3), (2,4) och (4,3) så det fungerar för mig! Jag kör på ditt förslag: C till (3,3), L till (2,4), P till (4,3). Jag börjar flytta mina objekt nu. Ska vi säga att vi meddelar varandra när vi är klara med alla flyttar?} \\
            }
        }
    }
    & & \\ \\

    \theutterance \stepcounter{utterance}  
    & & \multicolumn{4}{p{0.6\linewidth}}{
        \cellcolor[rgb]{0.9,0.9,0.9}{
            \makecell[{{p{\linewidth}}}]{
                \texttt{\tiny{[GM$\rangle$P2]}}
                \texttt{Ditt meddelande har vidarebefordrats till den andra spelaren.} \\
\\ 
\\ 
\texttt{Du spelar för närvarande runda 3 av maximalt 12.} \\
\texttt{Ni har tillsammans ackumulerat 1 av 8 straff.} \\
\texttt{Den andra spelaren skickade detta meddelande:} \\
\\ 
\\ 
\texttt{"Låt mig kolla min spelplan. Jag har tomma rutor på (3,3), (2,4) och (4,3) så det fungerar för mig! Jag kör på ditt förslag: C till (3,3), L till (2,4), P till (4,3). Jag börjar flytta mina objekt nu. Ska vi säga att vi meddelar varandra när vi är klara med alla flyttar?"} \\
\\ 
\\ 
\texttt{Vad är ditt nästa kommando?} \\
            }
        }
    }
    & \\ \\

    \theutterance \stepcounter{utterance}  
    & & & \multicolumn{4}{p{0.6\linewidth}}{
        \cellcolor[rgb]{0.9,0.9,0.9}{
            \makecell[{{p{\linewidth}}}]{
                \texttt{\tiny{[GM$\langle$P2]}}
                \texttt{Innan jag flyttar något måste jag kontrollera min spelplan för att se vilka rutor som är tomma.} \\
\\ 
\texttt{Min plan nu:} \\
\texttt{{-} C på (3,2)} \\
\texttt{{-} L på (1,7)} \\
\texttt{{-} P på (3,1)} \\
\\ 
\texttt{Målpositioner enligt överenskommelsen:} \\
\texttt{{-} C till (3,3)} \\
\texttt{{-} L till (2,4)} \\
\texttt{{-} P till (4,3)} \\
\\ 
\texttt{Jag kontrollerar (3,3): I min plan är (3,3) markerad med ◌, så den är tom.} \\
\texttt{Jag börjar med att flytta C från (3,2) till (3,3).} \\
\\ 
\texttt{FLYTTA: C, (3,3)} \\
            }
        }
    }
     \\ \\

    \theutterance \stepcounter{utterance}  
    & & & \multicolumn{2}{p{0.3\linewidth}}{
        \cellcolor[rgb]{0.9,0.9,0.9}{
            \makecell[{{p{\linewidth}}}]{
                \texttt{\tiny{[GM$|$GM]}}
                \texttt{Innan jag flyttar något måste jag kontrollera min spelplan för att se vilka rutor som är tomma.} \\
\\ 
\texttt{Min plan nu:} \\
\texttt{{-} C på (3,2)} \\
\texttt{{-} L på (1,7)} \\
\texttt{{-} P på (3,1)} \\
\\ 
\texttt{Målpositioner enligt överenskommelsen:} \\
\texttt{{-} C till (3,3)} \\
\texttt{{-} L till (2,4)} \\
\texttt{{-} P till (4,3)} \\
\\ 
\texttt{Jag kontrollerar (3,3): I min plan är (3,3) markerad med ◌, så den är tom.} \\
\texttt{Jag börjar med att flytta C från (3,2) till (3,3).} \\
\\ 
\texttt{FLYTTA: C, (3,3)} \\
            }
        }
    }
    & & \\ \\

    \theutterance \stepcounter{utterance}  
    & & & \multicolumn{2}{p{0.3\linewidth}}{
        \cellcolor[rgb]{0.9,0.9,0.9}{
            \makecell[{{p{\linewidth}}}]{
                \texttt{\tiny{[GM$|$GM]}}
                \texttt{Invalid format: head is not empty: 'Innan jag flyttar något måste jag kontrollera min spelplan för att se vilka rutor som är tomma.} \\
\\ 
\texttt{Min plan nu:} \\
\texttt{{-} C på (3,2)} \\
\texttt{{-} L på (1,7)} \\
\texttt{{-} P på (3,1)} \\
\\ 
\texttt{Målpositioner enligt överenskommelsen:} \\
\texttt{{-} C till (3,3)} \\
\texttt{{-} L till (2,4)} \\
\texttt{{-} P till (4,3)} \\
\\ 
\texttt{Jag kontrollerar (3,3): I min plan är (3,3) markerad med ◌, så den är tom.} \\
\texttt{Jag börjar med att flytta C från (3,2) till (3,3).} \\
\\ 
\texttt{'} \\
            }
        }
    }
    & & \\ \\

    \theutterance \stepcounter{utterance}  
    & & \multicolumn{4}{p{0.6\linewidth}}{
        \cellcolor[rgb]{0.9,0.9,0.9}{
            \makecell[{{p{\linewidth}}}]{
                \texttt{\tiny{[GM$\rangle$P2]}}
                \texttt{Straff: Ditt meddelande får inte innehålla något före kommandot!} \\
\\ 
\texttt{Se till att ditt svar endast innehåller antingen `SÄG: <MEDDELANDE>` eller `FLYTTA: <OBJEKT>, (<X>, <Y>)`, och inget annat!} \\
\\ 
\texttt{Ni har tillsammans ackumulerat 2 av 8 straff.} \\
\texttt{Vänligen försök igen!} \\
            }
        }
    }
    & \\ \\

    \theutterance \stepcounter{utterance}  
    & & & \multicolumn{4}{p{0.6\linewidth}}{
        \cellcolor[rgb]{0.9,0.9,0.9}{
            \makecell[{{p{\linewidth}}}]{
                \texttt{\tiny{[GM$\langle$P2]}}
                \texttt{FLYTTA: C, (3,3)} \\
            }
        }
    }
     \\ \\

    \theutterance \stepcounter{utterance}  
    & & & \multicolumn{2}{p{0.3\linewidth}}{
        \cellcolor[rgb]{0.9,0.9,0.9}{
            \makecell[{{p{\linewidth}}}]{
                \texttt{\tiny{[GM$|$GM]}}
                \texttt{FLYTTA: C, (3,3)} \\
            }
        }
    }
    & & \\ \\

    \theutterance \stepcounter{utterance}  
    & & & \multicolumn{2}{p{0.3\linewidth}}{
        \cellcolor[rgb]{0.9,0.9,0.9}{
            \makecell[{{p{\linewidth}}}]{
                \texttt{\tiny{[GM$|$GM]}}
                \texttt{C moved to (3, 3)} \\
            }
        }
    }
    & & \\ \\

    \theutterance \stepcounter{utterance}  
    & & \multicolumn{4}{p{0.6\linewidth}}{
        \cellcolor[rgb]{0.9,0.9,0.9}{
            \makecell[{{p{\linewidth}}}]{
                \texttt{\tiny{[P1$\langle$GM]}}
                \texttt{Ditt meddelande har vidarebefordrats till den andra spelaren.} \\
\\ 
\\ 
\texttt{Du spelar för närvarande runda 3 av maximalt 12.} \\
\texttt{Ni har tillsammans ackumulerat 2 av 8 straff.} \\
\texttt{Den andra spelaren flyttade ett objekt på sitt rutnät.} \\
\\ 
\\ 
\texttt{Vad är ditt nästa kommando?} \\
            }
        }
    }
    & \\ \\

    \theutterance \stepcounter{utterance}  
    & \multicolumn{4}{p{0.6\linewidth}}{
        \cellcolor[rgb]{0.9,0.9,0.9}{
            \makecell[{{p{\linewidth}}}]{
                \texttt{\tiny{[P1$\rangle$GM]}}
                \texttt{FLYTTA: C, (3,3)} \\
            }
        }
    }
    & & \\ \\

    \theutterance \stepcounter{utterance}  
    & & & \multicolumn{2}{p{0.3\linewidth}}{
        \cellcolor[rgb]{0.9,0.9,0.9}{
            \makecell[{{p{\linewidth}}}]{
                \texttt{\tiny{[GM$|$GM]}}
                \texttt{FLYTTA: C, (3,3)} \\
            }
        }
    }
    & & \\ \\

    \theutterance \stepcounter{utterance}  
    & & & \multicolumn{2}{p{0.3\linewidth}}{
        \cellcolor[rgb]{0.9,0.9,0.9}{
            \makecell[{{p{\linewidth}}}]{
                \texttt{\tiny{[GM$|$GM]}}
                \texttt{C moved to (3, 3)} \\
            }
        }
    }
    & & \\ \\

    \theutterance \stepcounter{utterance}  
    & & \multicolumn{4}{p{0.6\linewidth}}{
        \cellcolor[rgb]{0.9,0.9,0.9}{
            \makecell[{{p{\linewidth}}}]{
                \texttt{\tiny{[GM$\rangle$P2]}}
                \texttt{Flyttade 'C' till (3,3) framgångsrikt. Ditt uppdaterade rutnät ser ut så här:} \\
\texttt{```} \\
\texttt{ 1234567} \\
\texttt{╔════╤══╗} \\
\texttt{║◌◌P◌│◌◌║ 1} \\
\texttt{║◌◌◌◌├──╢ 2} \\
\texttt{║◌◌C◌│◌◌║ 3} \\
\texttt{║◌┌──┼──╢ 4} \\
\texttt{║◌│◌◌│◌◌║ 5} \\
\texttt{╟─┴┬─┤◌◌║ 6} \\
\texttt{║L◌│◌│◌◌║ 7} \\
\texttt{╚══╧═╧══╝} \\
\texttt{```} \\
\texttt{Du spelar för närvarande runda 4 av maximalt 12.} \\
\texttt{Ni har tillsammans ackumulerat 2 av 8 straff.} \\
\texttt{Den andra spelaren flyttade ett objekt på sitt rutnät.} \\
\\ 
\\ 
\texttt{Vad är ditt nästa kommando?} \\
            }
        }
    }
    & \\ \\

    \theutterance \stepcounter{utterance}  
    & & & \multicolumn{4}{p{0.6\linewidth}}{
        \cellcolor[rgb]{0.9,0.9,0.9}{
            \makecell[{{p{\linewidth}}}]{
                \texttt{\tiny{[GM$\langle$P2]}}
                \texttt{Nu ska jag flytta L från (1,7) till (2,4). Jag kontrollerar (2,4) på min plan {-} i rad 4, kolumn 2 ser jag en vägglinje, inte en tom ruta. Låt mig titta noggrannare.} \\
\\ 
\texttt{Rad 4: `║◌┌──┼──╢`} \\
\\ 
\texttt{Kolumn 2 i rad 4 är `┌` vilket är en väggdel, inte en tom ruta. (2,4) är alltså inte tillgänglig.} \\
\\ 
\texttt{Jag måste meddela den andra spelaren om detta problem.} \\
\\ 
\texttt{SÄG: Jag kan inte flytta L till (2,4) {-} den rutan innehåller en vägg på min plan. Kan vi hitta en annan position för L? Kanske (2,3) eller (2,5)? Vilka av dessa är lediga på din plan?} \\
            }
        }
    }
     \\ \\

    \theutterance \stepcounter{utterance}  
    & & & \multicolumn{2}{p{0.3\linewidth}}{
        \cellcolor[rgb]{0.9,0.9,0.9}{
            \makecell[{{p{\linewidth}}}]{
                \texttt{\tiny{[GM$|$GM]}}
                \texttt{Nu ska jag flytta L från (1,7) till (2,4). Jag kontrollerar (2,4) på min plan {-} i rad 4, kolumn 2 ser jag en vägglinje, inte en tom ruta. Låt mig titta noggrannare.} \\
\\ 
\texttt{Rad 4: `║◌┌──┼──╢`} \\
\\ 
\texttt{Kolumn 2 i rad 4 är `┌` vilket är en väggdel, inte en tom ruta. (2,4) är alltså inte tillgänglig.} \\
\\ 
\texttt{Jag måste meddela den andra spelaren om detta problem.} \\
\\ 
\texttt{SÄG: Jag kan inte flytta L till (2,4) {-} den rutan innehåller en vägg på min plan. Kan vi hitta en annan position för L? Kanske (2,3) eller (2,5)? Vilka av dessa är lediga på din plan?} \\
            }
        }
    }
    & & \\ \\

    \theutterance \stepcounter{utterance}  
    & & & \multicolumn{2}{p{0.3\linewidth}}{
        \cellcolor[rgb]{0.9,0.9,0.9}{
            \makecell[{{p{\linewidth}}}]{
                \texttt{\tiny{[GM$|$GM]}}
                \texttt{Invalid format: head is not empty: 'Nu ska jag flytta L från (1,7) till (2,4). Jag kontrollerar (2,4) på min plan {-} i rad 4, kolumn 2 ser jag en vägglinje, inte en tom ruta. Låt mig titta noggrannare.} \\
\\ 
\texttt{Rad 4: ║◌┌──┼──╢} \\
\\ 
\texttt{Kolumn 2 i rad 4 är ┌ vilket är en väggdel, inte en tom ruta. (2,4) är alltså inte tillgänglig.} \\
\\ 
\texttt{Jag måste meddela den andra spelaren om detta problem.} \\
\\ 
\texttt{'} \\
            }
        }
    }
    & & \\ \\

    \theutterance \stepcounter{utterance}  
    & & \multicolumn{4}{p{0.6\linewidth}}{
        \cellcolor[rgb]{0.9,0.9,0.9}{
            \makecell[{{p{\linewidth}}}]{
                \texttt{\tiny{[GM$\rangle$P2]}}
                \texttt{Straff: Ditt meddelande får inte innehålla något före kommandot!} \\
\\ 
\texttt{Se till att ditt svar endast innehåller antingen `SÄG: <MEDDELANDE>` eller `FLYTTA: <OBJEKT>, (<X>, <Y>)`, och inget annat!} \\
\\ 
\texttt{Ni har tillsammans ackumulerat 3 av 8 straff.} \\
\texttt{Vänligen försök igen!} \\
            }
        }
    }
    & \\ \\

    \theutterance \stepcounter{utterance}  
    & & & \multicolumn{4}{p{0.6\linewidth}}{
        \cellcolor[rgb]{0.9,0.9,0.9}{
            \makecell[{{p{\linewidth}}}]{
                \texttt{\tiny{[GM$\langle$P2]}}
                \texttt{SÄG: Jag kan inte flytta L till (2,4) {-} den rutan innehåller en vägg på min plan. Kan vi hitta en annan position för L? Kanske (2,3) eller (2,5)? Vilka av dessa är lediga på din plan?} \\
            }
        }
    }
     \\ \\

    \theutterance \stepcounter{utterance}  
    & & & \multicolumn{2}{p{0.3\linewidth}}{
        \cellcolor[rgb]{0.9,0.9,0.9}{
            \makecell[{{p{\linewidth}}}]{
                \texttt{\tiny{[GM$|$GM]}}
                \texttt{SÄG: Jag kan inte flytta L till (2,4) {-} den rutan innehåller en vägg på min plan. Kan vi hitta en annan position för L? Kanske (2,3) eller (2,5)? Vilka av dessa är lediga på din plan?} \\
            }
        }
    }
    & & \\ \\

    \theutterance \stepcounter{utterance}  
    & & \multicolumn{4}{p{0.6\linewidth}}{
        \cellcolor[rgb]{0.9,0.9,0.9}{
            \makecell[{{p{\linewidth}}}]{
                \texttt{\tiny{[P1$\langle$GM]}}
                \texttt{Flyttade 'C' till (3,3) framgångsrikt. Ditt uppdaterade rutnät ser ut så här:} \\
\texttt{```} \\
\texttt{ 1234567} \\
\texttt{╔════╤══╗} \\
\texttt{║◌◌◌◌│◌◌║ 1} \\
\texttt{║◌◌L◌├──╢ 2} \\
\texttt{║◌◌CP│◌◌║ 3} \\
\texttt{║◌┌──┼──╢ 4} \\
\texttt{║◌│◌◌│◌◌║ 5} \\
\texttt{╟─┴┬─┤◌◌║ 6} \\
\texttt{║◌◌│◌│◌◌║ 7} \\
\texttt{╚══╧═╧══╝} \\
\texttt{```} \\
\texttt{Du spelar för närvarande runda 4 av maximalt 12.} \\
\texttt{Ni har tillsammans ackumulerat 3 av 8 straff.} \\
\texttt{Den andra spelaren skickade detta meddelande:} \\
\\ 
\\ 
\texttt{"Jag kan inte flytta L till (2,4) {-} den rutan innehåller en vägg på min plan. Kan vi hitta en annan position för L? Kanske (2,3) eller (2,5)? Vilka av dessa är lediga på din plan?"} \\
\\ 
\\ 
\texttt{Vad är ditt nästa kommando?} \\
            }
        }
    }
    & \\ \\

    \theutterance \stepcounter{utterance}  
    & \multicolumn{4}{p{0.6\linewidth}}{
        \cellcolor[rgb]{0.9,0.9,0.9}{
            \makecell[{{p{\linewidth}}}]{
                \texttt{\tiny{[P1$\rangle$GM]}}
                \texttt{SÄG: På min plan är (2,3) upptagen av P just nu, men (2,5) är ledig. Ska vi sikta på L till (2,5) istället? Jag har C på (3,3) nu och P står kvar på (4,3).} \\
            }
        }
    }
    & & \\ \\

    \theutterance \stepcounter{utterance}  
    & & & \multicolumn{2}{p{0.3\linewidth}}{
        \cellcolor[rgb]{0.9,0.9,0.9}{
            \makecell[{{p{\linewidth}}}]{
                \texttt{\tiny{[GM$|$GM]}}
                \texttt{SÄG: På min plan är (2,3) upptagen av P just nu, men (2,5) är ledig. Ska vi sikta på L till (2,5) istället? Jag har C på (3,3) nu och P står kvar på (4,3).} \\
            }
        }
    }
    & & \\ \\

    \theutterance \stepcounter{utterance}  
    & & \multicolumn{4}{p{0.6\linewidth}}{
        \cellcolor[rgb]{0.9,0.9,0.9}{
            \makecell[{{p{\linewidth}}}]{
                \texttt{\tiny{[GM$\rangle$P2]}}
                \texttt{Ditt meddelande har vidarebefordrats till den andra spelaren.} \\
\\ 
\\ 
\texttt{Du spelar för närvarande runda 5 av maximalt 12.} \\
\texttt{Ni har tillsammans ackumulerat 3 av 8 straff.} \\
\texttt{Den andra spelaren skickade detta meddelande:} \\
\\ 
\\ 
\texttt{"På min plan är (2,3) upptagen av P just nu, men (2,5) är ledig. Ska vi sikta på L till (2,5) istället? Jag har C på (3,3) nu och P står kvar på (4,3)."} \\
\\ 
\\ 
\texttt{Vad är ditt nästa kommando?} \\
            }
        }
    }
    & \\ \\

    \theutterance \stepcounter{utterance}  
    & & & \multicolumn{4}{p{0.6\linewidth}}{
        \cellcolor[rgb]{0.9,0.9,0.9}{
            \makecell[{{p{\linewidth}}}]{
                \texttt{\tiny{[GM$\langle$P2]}}
                \texttt{Jag kontrollerar (2,5) på min plan. Rad 5: `║◌│◌◌│◌◌║`} \\
\\ 
\texttt{Kolumn 2 i rad 5 är `│` vilket är en vägg, inte en tom ruta. (2,5) är alltså inte tillgänglig heller.} \\
\\ 
\texttt{SÄG: (2,5) är också en vägg på min plan. Mina väggar är besvärliga! Kan vi försöka med (4,5) eller (5,5) för L? Eller har du andra lediga rutor i mitten av planen?} \\
            }
        }
    }
     \\ \\

    \theutterance \stepcounter{utterance}  
    & & & \multicolumn{2}{p{0.3\linewidth}}{
        \cellcolor[rgb]{0.9,0.9,0.9}{
            \makecell[{{p{\linewidth}}}]{
                \texttt{\tiny{[GM$|$GM]}}
                \texttt{Jag kontrollerar (2,5) på min plan. Rad 5: `║◌│◌◌│◌◌║`} \\
\\ 
\texttt{Kolumn 2 i rad 5 är `│` vilket är en vägg, inte en tom ruta. (2,5) är alltså inte tillgänglig heller.} \\
\\ 
\texttt{SÄG: (2,5) är också en vägg på min plan. Mina väggar är besvärliga! Kan vi försöka med (4,5) eller (5,5) för L? Eller har du andra lediga rutor i mitten av planen?} \\
            }
        }
    }
    & & \\ \\

    \theutterance \stepcounter{utterance}  
    & & & \multicolumn{2}{p{0.3\linewidth}}{
        \cellcolor[rgb]{0.9,0.9,0.9}{
            \makecell[{{p{\linewidth}}}]{
                \texttt{\tiny{[GM$|$GM]}}
                \texttt{Invalid format: head is not empty: 'Jag kontrollerar (2,5) på min plan. Rad 5: ║◌│◌◌│◌◌║} \\
\\ 
\texttt{Kolumn 2 i rad 5 är │ vilket är en vägg, inte en tom ruta. (2,5) är alltså inte tillgänglig heller.} \\
\\ 
\texttt{'} \\
            }
        }
    }
    & & \\ \\

    \theutterance \stepcounter{utterance}  
    & & \multicolumn{4}{p{0.6\linewidth}}{
        \cellcolor[rgb]{0.9,0.9,0.9}{
            \makecell[{{p{\linewidth}}}]{
                \texttt{\tiny{[GM$\rangle$P2]}}
                \texttt{Straff: Ditt meddelande får inte innehålla något före kommandot!} \\
\\ 
\texttt{Se till att ditt svar endast innehåller antingen `SÄG: <MEDDELANDE>` eller `FLYTTA: <OBJEKT>, (<X>, <Y>)`, och inget annat!} \\
\\ 
\texttt{Ni har tillsammans ackumulerat 4 av 8 straff.} \\
\texttt{Vänligen försök igen!} \\
            }
        }
    }
    & \\ \\

    \theutterance \stepcounter{utterance}  
    & & & \multicolumn{4}{p{0.6\linewidth}}{
        \cellcolor[rgb]{0.9,0.9,0.9}{
            \makecell[{{p{\linewidth}}}]{
                \texttt{\tiny{[GM$\langle$P2]}}
                \texttt{SÄG: (2,5) är också en vägg på min plan. Mina väggar är besvärliga! Kan vi försöka med (4,5) eller (5,5) för L? Eller har du andra lediga rutor i mitten av planen?} \\
            }
        }
    }
     \\ \\

    \theutterance \stepcounter{utterance}  
    & & & \multicolumn{2}{p{0.3\linewidth}}{
        \cellcolor[rgb]{0.9,0.9,0.9}{
            \makecell[{{p{\linewidth}}}]{
                \texttt{\tiny{[GM$|$GM]}}
                \texttt{SÄG: (2,5) är också en vägg på min plan. Mina väggar är besvärliga! Kan vi försöka med (4,5) eller (5,5) för L? Eller har du andra lediga rutor i mitten av planen?} \\
            }
        }
    }
    & & \\ \\

    \theutterance \stepcounter{utterance}  
    & & \multicolumn{4}{p{0.6\linewidth}}{
        \cellcolor[rgb]{0.9,0.9,0.9}{
            \makecell[{{p{\linewidth}}}]{
                \texttt{\tiny{[P1$\langle$GM]}}
                \texttt{Ditt meddelande har vidarebefordrats till den andra spelaren.} \\
\\ 
\\ 
\texttt{Du spelar för närvarande runda 5 av maximalt 12.} \\
\texttt{Ni har tillsammans ackumulerat 4 av 8 straff.} \\
\texttt{Den andra spelaren skickade detta meddelande:} \\
\\ 
\\ 
\texttt{"(2,5) är också en vägg på min plan. Mina väggar är besvärliga! Kan vi försöka med (4,5) eller (5,5) för L? Eller har du andra lediga rutor i mitten av planen?"} \\
\\ 
\\ 
\texttt{Vad är ditt nästa kommando?} \\
            }
        }
    }
    & \\ \\

    \theutterance \stepcounter{utterance}  
    & \multicolumn{4}{p{0.6\linewidth}}{
        \cellcolor[rgb]{0.9,0.9,0.9}{
            \makecell[{{p{\linewidth}}}]{
                \texttt{\tiny{[P1$\rangle$GM]}}
                \texttt{SÄG: Låt mig kolla. (4,5) och (5,5) är båda lediga på min plan. (5,3) är också ledig. Vilken av dessa fungerar för dig? Vi kan också titta på (6,3), (6,4), (6,5) eller (6,6) om någon av dem passar bättre med dina väggar.} \\
            }
        }
    }
    & & \\ \\

    \theutterance \stepcounter{utterance}  
    & & & \multicolumn{2}{p{0.3\linewidth}}{
        \cellcolor[rgb]{0.9,0.9,0.9}{
            \makecell[{{p{\linewidth}}}]{
                \texttt{\tiny{[GM$|$GM]}}
                \texttt{SÄG: Låt mig kolla. (4,5) och (5,5) är båda lediga på min plan. (5,3) är också ledig. Vilken av dessa fungerar för dig? Vi kan också titta på (6,3), (6,4), (6,5) eller (6,6) om någon av dem passar bättre med dina väggar.} \\
            }
        }
    }
    & & \\ \\

    \theutterance \stepcounter{utterance}  
    & & \multicolumn{4}{p{0.6\linewidth}}{
        \cellcolor[rgb]{0.9,0.9,0.9}{
            \makecell[{{p{\linewidth}}}]{
                \texttt{\tiny{[GM$\rangle$P2]}}
                \texttt{Ditt meddelande har vidarebefordrats till den andra spelaren.} \\
\\ 
\\ 
\texttt{Du spelar för närvarande runda 6 av maximalt 12.} \\
\texttt{Ni har tillsammans ackumulerat 4 av 8 straff.} \\
\texttt{Den andra spelaren skickade detta meddelande:} \\
\\ 
\\ 
\texttt{"Låt mig kolla. (4,5) och (5,5) är båda lediga på min plan. (5,3) är också ledig. Vilken av dessa fungerar för dig? Vi kan också titta på (6,3), (6,4), (6,5) eller (6,6) om någon av dem passar bättre med dina väggar."} \\
\\ 
\\ 
\texttt{Vad är ditt nästa kommando?} \\
            }
        }
    }
    & \\ \\

    \theutterance \stepcounter{utterance}  
    & & & \multicolumn{4}{p{0.6\linewidth}}{
        \cellcolor[rgb]{0.9,0.9,0.9}{
            \makecell[{{p{\linewidth}}}]{
                \texttt{\tiny{[GM$\langle$P2]}}
                \texttt{Jag kontrollerar dessa rutor på min plan:} \\
\\ 
\texttt{{-} (4,5): Rad 5, kolumn 4 är `◌` {-} tom!} \\
\texttt{{-} (5,5): Rad 5, kolumn 5 är `│` {-} vägg.} \\
\texttt{{-} (5,3): Rad 3, kolumn 5 är `│` {-} vägg.} \\
\texttt{{-} (6,3): Rad 3, kolumn 6 är `◌` {-} tom!} \\
\texttt{{-} (6,4): Rad 4, kolumn 6 är `─` {-} vägg.} \\
\texttt{{-} (6,5): Rad 5, kolumn 6 är `◌` {-} tom!} \\
\texttt{{-} (6,6): Rad 6, kolumn 6 är `◌` {-} tom!} \\
\\ 
\texttt{(4,5) fungerar utmärkt. Säg att vi kör på den.} \\
\\ 
\texttt{SÄG: (4,5) fungerar för mig! Vi kör L till (4,5). Då har vi C på (3,3), L på (4,5), P på (4,3). Okej?} \\
            }
        }
    }
     \\ \\

    \theutterance \stepcounter{utterance}  
    & & & \multicolumn{2}{p{0.3\linewidth}}{
        \cellcolor[rgb]{0.9,0.9,0.9}{
            \makecell[{{p{\linewidth}}}]{
                \texttt{\tiny{[GM$|$GM]}}
                \texttt{Jag kontrollerar dessa rutor på min plan:} \\
\\ 
\texttt{{-} (4,5): Rad 5, kolumn 4 är `◌` {-} tom!} \\
\texttt{{-} (5,5): Rad 5, kolumn 5 är `│` {-} vägg.} \\
\texttt{{-} (5,3): Rad 3, kolumn 5 är `│` {-} vägg.} \\
\texttt{{-} (6,3): Rad 3, kolumn 6 är `◌` {-} tom!} \\
\texttt{{-} (6,4): Rad 4, kolumn 6 är `─` {-} vägg.} \\
\texttt{{-} (6,5): Rad 5, kolumn 6 är `◌` {-} tom!} \\
\texttt{{-} (6,6): Rad 6, kolumn 6 är `◌` {-} tom!} \\
\\ 
\texttt{(4,5) fungerar utmärkt. Säg att vi kör på den.} \\
\\ 
\texttt{SÄG: (4,5) fungerar för mig! Vi kör L till (4,5). Då har vi C på (3,3), L på (4,5), P på (4,3). Okej?} \\
            }
        }
    }
    & & \\ \\

    \theutterance \stepcounter{utterance}  
    & & & \multicolumn{2}{p{0.3\linewidth}}{
        \cellcolor[rgb]{0.9,0.9,0.9}{
            \makecell[{{p{\linewidth}}}]{
                \texttt{\tiny{[GM$|$GM]}}
                \texttt{Invalid format: head is not empty: 'Jag kontrollerar dessa rutor på min plan:} \\
\\ 
\texttt{{-} (4,5): Rad 5, kolumn 4 är ◌ {-} tom!} \\
\texttt{{-} (5,5): Rad 5, kolumn 5 är │ {-} vägg.} \\
\texttt{{-} (5,3): Rad 3, kolumn 5 är │ {-} vägg.} \\
\texttt{{-} (6,3): Rad 3, kolumn 6 är ◌ {-} tom!} \\
\texttt{{-} (6,4): Rad 4, kolumn 6 är ─ {-} vägg.} \\
\texttt{{-} (6,5): Rad 5, kolumn 6 är ◌ {-} tom!} \\
\texttt{{-} (6,6): Rad 6, kolumn 6 är ◌ {-} tom!} \\
\\ 
\texttt{(4,5) fungerar utmärkt. Säg att vi kör på den.} \\
\\ 
\texttt{'} \\
            }
        }
    }
    & & \\ \\

    \theutterance \stepcounter{utterance}  
    & & \multicolumn{4}{p{0.6\linewidth}}{
        \cellcolor[rgb]{0.9,0.9,0.9}{
            \makecell[{{p{\linewidth}}}]{
                \texttt{\tiny{[GM$\rangle$P2]}}
                \texttt{Straff: Ditt meddelande får inte innehålla något före kommandot!} \\
\\ 
\texttt{Se till att ditt svar endast innehåller antingen `SÄG: <MEDDELANDE>` eller `FLYTTA: <OBJEKT>, (<X>, <Y>)`, och inget annat!} \\
\\ 
\texttt{Ni har tillsammans ackumulerat 5 av 8 straff.} \\
\texttt{Vänligen försök igen!} \\
            }
        }
    }
    & \\ \\

    \theutterance \stepcounter{utterance}  
    & & & \multicolumn{4}{p{0.6\linewidth}}{
        \cellcolor[rgb]{0.9,0.9,0.9}{
            \makecell[{{p{\linewidth}}}]{
                \texttt{\tiny{[GM$\langle$P2]}}
                \texttt{SÄG: (4,5) fungerar för mig! Vi kör L till (4,5). Då har vi C på (3,3), L på (4,5), P på (4,3). Okej?} \\
            }
        }
    }
     \\ \\

    \theutterance \stepcounter{utterance}  
    & & & \multicolumn{2}{p{0.3\linewidth}}{
        \cellcolor[rgb]{0.9,0.9,0.9}{
            \makecell[{{p{\linewidth}}}]{
                \texttt{\tiny{[GM$|$GM]}}
                \texttt{SÄG: (4,5) fungerar för mig! Vi kör L till (4,5). Då har vi C på (3,3), L på (4,5), P på (4,3). Okej?} \\
            }
        }
    }
    & & \\ \\

    \theutterance \stepcounter{utterance}  
    & & \multicolumn{4}{p{0.6\linewidth}}{
        \cellcolor[rgb]{0.9,0.9,0.9}{
            \makecell[{{p{\linewidth}}}]{
                \texttt{\tiny{[P1$\langle$GM]}}
                \texttt{Ditt meddelande har vidarebefordrats till den andra spelaren.} \\
\\ 
\\ 
\texttt{Du spelar för närvarande runda 6 av maximalt 12.} \\
\texttt{Ni har tillsammans ackumulerat 5 av 8 straff.} \\
\texttt{Den andra spelaren skickade detta meddelande:} \\
\\ 
\\ 
\texttt{"(4,5) fungerar för mig! Vi kör L till (4,5). Då har vi C på (3,3), L på (4,5), P på (4,3). Okej?"} \\
\\ 
\\ 
\texttt{Vad är ditt nästa kommando?} \\
            }
        }
    }
    & \\ \\

    \theutterance \stepcounter{utterance}  
    & \multicolumn{4}{p{0.6\linewidth}}{
        \cellcolor[rgb]{0.9,0.9,0.9}{
            \makecell[{{p{\linewidth}}}]{
                \texttt{\tiny{[P1$\rangle$GM]}}
                \texttt{SÄG: Perfekt! Jag bekräftar: C på (3,3), L på (4,5), P på (4,3). Jag flyttar L och P nu.} \\
            }
        }
    }
    & & \\ \\

    \theutterance \stepcounter{utterance}  
    & & & \multicolumn{2}{p{0.3\linewidth}}{
        \cellcolor[rgb]{0.9,0.9,0.9}{
            \makecell[{{p{\linewidth}}}]{
                \texttt{\tiny{[GM$|$GM]}}
                \texttt{SÄG: Perfekt! Jag bekräftar: C på (3,3), L på (4,5), P på (4,3). Jag flyttar L och P nu.} \\
            }
        }
    }
    & & \\ \\

    \theutterance \stepcounter{utterance}  
    & & \multicolumn{4}{p{0.6\linewidth}}{
        \cellcolor[rgb]{0.9,0.9,0.9}{
            \makecell[{{p{\linewidth}}}]{
                \texttt{\tiny{[GM$\rangle$P2]}}
                \texttt{Ditt meddelande har vidarebefordrats till den andra spelaren.} \\
\\ 
\\ 
\texttt{Du spelar för närvarande runda 7 av maximalt 12.} \\
\texttt{Ni har tillsammans ackumulerat 5 av 8 straff.} \\
\texttt{Den andra spelaren skickade detta meddelande:} \\
\\ 
\\ 
\texttt{"Perfekt! Jag bekräftar: C på (3,3), L på (4,5), P på (4,3). Jag flyttar L och P nu."} \\
\\ 
\\ 
\texttt{Vad är ditt nästa kommando?} \\
            }
        }
    }
    & \\ \\

    \theutterance \stepcounter{utterance}  
    & & & \multicolumn{4}{p{0.6\linewidth}}{
        \cellcolor[rgb]{0.9,0.9,0.9}{
            \makecell[{{p{\linewidth}}}]{
                \texttt{\tiny{[GM$\langle$P2]}}
                \texttt{Nu flyttar jag L från (1,7) till (4,5) och P från (3,1) till (4,3). Börjar med L.} \\
\\ 
\texttt{FLYTTA: L, (4,5)} \\
            }
        }
    }
     \\ \\

    \theutterance \stepcounter{utterance}  
    & & & \multicolumn{2}{p{0.3\linewidth}}{
        \cellcolor[rgb]{0.9,0.9,0.9}{
            \makecell[{{p{\linewidth}}}]{
                \texttt{\tiny{[GM$|$GM]}}
                \texttt{Nu flyttar jag L från (1,7) till (4,5) och P från (3,1) till (4,3). Börjar med L.} \\
\\ 
\texttt{FLYTTA: L, (4,5)} \\
            }
        }
    }
    & & \\ \\

    \theutterance \stepcounter{utterance}  
    & & & \multicolumn{2}{p{0.3\linewidth}}{
        \cellcolor[rgb]{0.9,0.9,0.9}{
            \makecell[{{p{\linewidth}}}]{
                \texttt{\tiny{[GM$|$GM]}}
                \texttt{Invalid format: head is not empty: 'Nu flyttar jag L från (1,7) till (4,5) och P från (3,1) till (4,3). Börjar med L.} \\
\\ 
\texttt{'} \\
            }
        }
    }
    & & \\ \\

    \theutterance \stepcounter{utterance}  
    & & \multicolumn{4}{p{0.6\linewidth}}{
        \cellcolor[rgb]{0.9,0.9,0.9}{
            \makecell[{{p{\linewidth}}}]{
                \texttt{\tiny{[GM$\rangle$P2]}}
                \texttt{Straff: Ditt meddelande får inte innehålla något före kommandot!} \\
\\ 
\texttt{Se till att ditt svar endast innehåller antingen `SÄG: <MEDDELANDE>` eller `FLYTTA: <OBJEKT>, (<X>, <Y>)`, och inget annat!} \\
\\ 
\texttt{Ni har tillsammans ackumulerat 6 av 8 straff.} \\
\texttt{Vänligen försök igen!} \\
            }
        }
    }
    & \\ \\

    \theutterance \stepcounter{utterance}  
    & & & \multicolumn{4}{p{0.6\linewidth}}{
        \cellcolor[rgb]{0.9,0.9,0.9}{
            \makecell[{{p{\linewidth}}}]{
                \texttt{\tiny{[GM$\langle$P2]}}
                \texttt{FLYTTA: L, (4,5)} \\
            }
        }
    }
     \\ \\

    \theutterance \stepcounter{utterance}  
    & & & \multicolumn{2}{p{0.3\linewidth}}{
        \cellcolor[rgb]{0.9,0.9,0.9}{
            \makecell[{{p{\linewidth}}}]{
                \texttt{\tiny{[GM$|$GM]}}
                \texttt{FLYTTA: L, (4,5)} \\
            }
        }
    }
    & & \\ \\

    \theutterance \stepcounter{utterance}  
    & & & \multicolumn{2}{p{0.3\linewidth}}{
        \cellcolor[rgb]{0.9,0.9,0.9}{
            \makecell[{{p{\linewidth}}}]{
                \texttt{\tiny{[GM$|$GM]}}
                \texttt{L moved to (4, 5)} \\
            }
        }
    }
    & & \\ \\

    \theutterance \stepcounter{utterance}  
    & & \multicolumn{4}{p{0.6\linewidth}}{
        \cellcolor[rgb]{0.9,0.9,0.9}{
            \makecell[{{p{\linewidth}}}]{
                \texttt{\tiny{[P1$\langle$GM]}}
                \texttt{Ditt meddelande har vidarebefordrats till den andra spelaren.} \\
\\ 
\\ 
\texttt{Du spelar för närvarande runda 7 av maximalt 12.} \\
\texttt{Ni har tillsammans ackumulerat 6 av 8 straff.} \\
\texttt{Den andra spelaren flyttade ett objekt på sitt rutnät.} \\
\\ 
\\ 
\texttt{Vad är ditt nästa kommando?} \\
            }
        }
    }
    & \\ \\

    \theutterance \stepcounter{utterance}  
    & \multicolumn{4}{p{0.6\linewidth}}{
        \cellcolor[rgb]{0.9,0.9,0.9}{
            \makecell[{{p{\linewidth}}}]{
                \texttt{\tiny{[P1$\rangle$GM]}}
                \texttt{FLYTTA: L, (4,5)} \\
            }
        }
    }
    & & \\ \\

    \theutterance \stepcounter{utterance}  
    & & & \multicolumn{2}{p{0.3\linewidth}}{
        \cellcolor[rgb]{0.9,0.9,0.9}{
            \makecell[{{p{\linewidth}}}]{
                \texttt{\tiny{[GM$|$GM]}}
                \texttt{FLYTTA: L, (4,5)} \\
            }
        }
    }
    & & \\ \\

    \theutterance \stepcounter{utterance}  
    & & & \multicolumn{2}{p{0.3\linewidth}}{
        \cellcolor[rgb]{0.9,0.9,0.9}{
            \makecell[{{p{\linewidth}}}]{
                \texttt{\tiny{[GM$|$GM]}}
                \texttt{L moved to (4, 5)} \\
            }
        }
    }
    & & \\ \\

    \theutterance \stepcounter{utterance}  
    & & \multicolumn{4}{p{0.6\linewidth}}{
        \cellcolor[rgb]{0.9,0.9,0.9}{
            \makecell[{{p{\linewidth}}}]{
                \texttt{\tiny{[GM$\rangle$P2]}}
                \texttt{Flyttade 'L' till (4,5) framgångsrikt. Ditt uppdaterade rutnät ser ut så här:} \\
\texttt{```} \\
\texttt{ 1234567} \\
\texttt{╔════╤══╗} \\
\texttt{║◌◌P◌│◌◌║ 1} \\
\texttt{║◌◌◌◌├──╢ 2} \\
\texttt{║◌◌C◌│◌◌║ 3} \\
\texttt{║◌┌──┼──╢ 4} \\
\texttt{║◌│◌L│◌◌║ 5} \\
\texttt{╟─┴┬─┤◌◌║ 6} \\
\texttt{║◌◌│◌│◌◌║ 7} \\
\texttt{╚══╧═╧══╝} \\
\texttt{```} \\
\texttt{Du spelar för närvarande runda 8 av maximalt 12.} \\
\texttt{Ni har tillsammans ackumulerat 6 av 8 straff.} \\
\texttt{Den andra spelaren flyttade ett objekt på sitt rutnät.} \\
\\ 
\\ 
\texttt{Vad är ditt nästa kommando?} \\
            }
        }
    }
    & \\ \\

    \theutterance \stepcounter{utterance}  
    & & & \multicolumn{4}{p{0.6\linewidth}}{
        \cellcolor[rgb]{0.9,0.9,0.9}{
            \makecell[{{p{\linewidth}}}]{
                \texttt{\tiny{[GM$\langle$P2]}}
                \texttt{Nu flyttar jag P från (3,1) till (4,3). Kontrollerar (4,3) på min plan: rad 3, kolumn 4 är `◌` {-} tom.} \\
\\ 
\texttt{FLYTTA: P, (4,3)} \\
            }
        }
    }
     \\ \\

    \theutterance \stepcounter{utterance}  
    & & & \multicolumn{2}{p{0.3\linewidth}}{
        \cellcolor[rgb]{0.9,0.9,0.9}{
            \makecell[{{p{\linewidth}}}]{
                \texttt{\tiny{[GM$|$GM]}}
                \texttt{Nu flyttar jag P från (3,1) till (4,3). Kontrollerar (4,3) på min plan: rad 3, kolumn 4 är `◌` {-} tom.} \\
\\ 
\texttt{FLYTTA: P, (4,3)} \\
            }
        }
    }
    & & \\ \\

    \theutterance \stepcounter{utterance}  
    & & & \multicolumn{2}{p{0.3\linewidth}}{
        \cellcolor[rgb]{0.9,0.9,0.9}{
            \makecell[{{p{\linewidth}}}]{
                \texttt{\tiny{[GM$|$GM]}}
                \texttt{Invalid format: head is not empty: 'Nu flyttar jag P från (3,1) till (4,3). Kontrollerar (4,3) på min plan: rad 3, kolumn 4 är ◌ {-} tom.} \\
\\ 
\texttt{'} \\
            }
        }
    }
    & & \\ \\

    \theutterance \stepcounter{utterance}  
    & & \multicolumn{4}{p{0.6\linewidth}}{
        \cellcolor[rgb]{0.9,0.9,0.9}{
            \makecell[{{p{\linewidth}}}]{
                \texttt{\tiny{[GM$\rangle$P2]}}
                \texttt{Straff: Ditt meddelande får inte innehålla något före kommandot!} \\
\\ 
\texttt{Se till att ditt svar endast innehåller antingen `SÄG: <MEDDELANDE>` eller `FLYTTA: <OBJEKT>, (<X>, <Y>)`, och inget annat!} \\
\\ 
\texttt{Ni har tillsammans ackumulerat 7 av 8 straff.} \\
\texttt{Vänligen försök igen!} \\
            }
        }
    }
    & \\ \\

    \theutterance \stepcounter{utterance}  
    & & & \multicolumn{4}{p{0.6\linewidth}}{
        \cellcolor[rgb]{0.9,0.9,0.9}{
            \makecell[{{p{\linewidth}}}]{
                \texttt{\tiny{[GM$\langle$P2]}}
                \texttt{FLYTTA: P, (4,3)} \\
            }
        }
    }
     \\ \\

    \theutterance \stepcounter{utterance}  
    & & & \multicolumn{2}{p{0.3\linewidth}}{
        \cellcolor[rgb]{0.9,0.9,0.9}{
            \makecell[{{p{\linewidth}}}]{
                \texttt{\tiny{[GM$|$GM]}}
                \texttt{FLYTTA: P, (4,3)} \\
            }
        }
    }
    & & \\ \\

    \theutterance \stepcounter{utterance}  
    & & & \multicolumn{2}{p{0.3\linewidth}}{
        \cellcolor[rgb]{0.9,0.9,0.9}{
            \makecell[{{p{\linewidth}}}]{
                \texttt{\tiny{[GM$|$GM]}}
                \texttt{P moved to (4, 3)} \\
            }
        }
    }
    & & \\ \\

    \theutterance \stepcounter{utterance}  
    & & \multicolumn{4}{p{0.6\linewidth}}{
        \cellcolor[rgb]{0.9,0.9,0.9}{
            \makecell[{{p{\linewidth}}}]{
                \texttt{\tiny{[P1$\langle$GM]}}
                \texttt{Flyttade 'L' till (4,5) framgångsrikt. Ditt uppdaterade rutnät ser ut så här:} \\
\texttt{```} \\
\texttt{ 1234567} \\
\texttt{╔════╤══╗} \\
\texttt{║◌◌◌◌│◌◌║ 1} \\
\texttt{║◌◌◌◌├──╢ 2} \\
\texttt{║◌◌CP│◌◌║ 3} \\
\texttt{║◌┌──┼──╢ 4} \\
\texttt{║◌│◌L│◌◌║ 5} \\
\texttt{╟─┴┬─┤◌◌║ 6} \\
\texttt{║◌◌│◌│◌◌║ 7} \\
\texttt{╚══╧═╧══╝} \\
\texttt{```} \\
\texttt{Du spelar för närvarande runda 8 av maximalt 12.} \\
\texttt{Ni har tillsammans ackumulerat 7 av 8 straff.} \\
\texttt{Den andra spelaren flyttade ett objekt på sitt rutnät.} \\
\\ 
\\ 
\texttt{Vad är ditt nästa kommando?} \\
            }
        }
    }
    & \\ \\

    \theutterance \stepcounter{utterance}  
    & \multicolumn{4}{p{0.6\linewidth}}{
        \cellcolor[rgb]{0.9,0.9,0.9}{
            \makecell[{{p{\linewidth}}}]{
                \texttt{\tiny{[P1$\rangle$GM]}}
                \texttt{FLYTTA: P, (4,3)} \\
            }
        }
    }
    & & \\ \\

    \theutterance \stepcounter{utterance}  
    & & & \multicolumn{2}{p{0.3\linewidth}}{
        \cellcolor[rgb]{0.9,0.9,0.9}{
            \makecell[{{p{\linewidth}}}]{
                \texttt{\tiny{[GM$|$GM]}}
                \texttt{FLYTTA: P, (4,3)} \\
            }
        }
    }
    & & \\ \\

    \theutterance \stepcounter{utterance}  
    & & & \multicolumn{2}{p{0.3\linewidth}}{
        \cellcolor[rgb]{0.9,0.9,0.9}{
            \makecell[{{p{\linewidth}}}]{
                \texttt{\tiny{[GM$|$GM]}}
                \texttt{Straff: (4,3) är inte tomt, utan innehåller 'P'.} \\
\\ 
\texttt{Ni har tillsammans ackumulerat 8 av 8 straff.} \\
\texttt{Vänligen försök igen!} \\
            }
        }
    }
    & & \\ \\

    \theutterance \stepcounter{utterance}  
    & & \multicolumn{4}{p{0.6\linewidth}}{
        \cellcolor[rgb]{0.9,0.9,0.9}{
            \makecell[{{p{\linewidth}}}]{
                \texttt{\tiny{[P1$\langle$GM]}}
                \texttt{Straff: (4,3) är inte tomt, utan innehåller 'P'.} \\
\\ 
\texttt{Ni har tillsammans ackumulerat 8 av 8 straff.} \\
\texttt{Vänligen försök igen!} \\
            }
        }
    }
    & \\ \\

    \theutterance \stepcounter{utterance}  
    & \multicolumn{4}{p{0.6\linewidth}}{
        \cellcolor[rgb]{0.9,0.9,0.9}{
            \makecell[{{p{\linewidth}}}]{
                \texttt{\tiny{[P1$\rangle$GM]}}
                \texttt{Åh nej, P står redan på (4,3)! Det är ju redan rätt position. Låt mig kolla min plan igen. C är på (3,3), P är redan på (4,3), och L är på (4,5). Alla är på rätt plats!} \\
\\ 
\texttt{SÄG: Jag är klar! Alla mina objekt står rätt: C på (3,3), P på (4,3), L på (4,5). Är du också klar? Ska vi avsluta? färdig?} \\
            }
        }
    }
    & & \\ \\

    \theutterance \stepcounter{utterance}  
    & & & \multicolumn{2}{p{0.3\linewidth}}{
        \cellcolor[rgb]{0.9,0.9,0.9}{
            \makecell[{{p{\linewidth}}}]{
                \texttt{\tiny{[GM$|$GM]}}
                \texttt{Åh nej, P står redan på (4,3)! Det är ju redan rätt position. Låt mig kolla min plan igen. C är på (3,3), P är redan på (4,3), och L är på (4,5). Alla är på rätt plats!} \\
\\ 
\texttt{SÄG: Jag är klar! Alla mina objekt står rätt: C på (3,3), P på (4,3), L på (4,5). Är du också klar? Ska vi avsluta? färdig?} \\
            }
        }
    }
    & & \\ \\

    \theutterance \stepcounter{utterance}  
    & & & \multicolumn{2}{p{0.3\linewidth}}{
        \cellcolor[rgb]{0.9,0.9,0.9}{
            \makecell[{{p{\linewidth}}}]{
                \texttt{\tiny{[GM$|$GM]}}
                \texttt{Invalid format: head is not empty: 'Åh nej, P står redan på (4,3)! Det är ju redan rätt position. Låt mig kolla min plan igen. C är på (3,3), P är redan på (4,3), och L är på (4,5). Alla är på rätt plats!} \\
\\ 
\texttt{'} \\
            }
        }
    }
    & & \\ \\

    \theutterance \stepcounter{utterance}  
    & & & \multicolumn{2}{p{0.3\linewidth}}{
        \cellcolor[rgb]{0.9,0.9,0.9}{
            \makecell[{{p{\linewidth}}}]{
                \texttt{\tiny{[GM$|$GM]}}
                \texttt{Maximum number of penalties exceeded} \\
            }
        }
    }
    & & \\ \\

    \theutterance \stepcounter{utterance}  
    & & & \multicolumn{2}{p{0.3\linewidth}}{
        \cellcolor[rgb]{0.9,0.9,0.9}{
            \makecell[{{p{\linewidth}}}]{
                \texttt{\tiny{[GM$|$GM]}}
                \texttt{Initial states:} \\
\texttt{```} \\
\texttt{Player 1:} \\
\texttt{ 1234567} \\
\texttt{╔════╤══╗} \\
\texttt{║◌C◌◌│◌◌║ 1} \\
\texttt{║◌◌L◌├──╢ 2} \\
\texttt{║◌◌◌P│◌◌║ 3} \\
\texttt{║◌┌──┼──╢ 4} \\
\texttt{║◌│◌◌│◌◌║ 5} \\
\texttt{╟─┴┬─┤◌◌║ 6} \\
\texttt{║◌◌│◌│◌◌║ 7} \\
\texttt{╚══╧═╧══╝} \\
\\ 
\\ 
\texttt{Player 2:} \\
\texttt{ 1234567} \\
\texttt{╔════╤══╗} \\
\texttt{║◌◌P◌│◌◌║ 1} \\
\texttt{║◌◌C◌├──╢ 2} \\
\texttt{║◌◌◌◌│◌◌║ 3} \\
\texttt{║◌┌──┼──╢ 4} \\
\texttt{║◌│◌◌│◌◌║ 5} \\
\texttt{╟─┴┬─┤◌◌║ 6} \\
\texttt{║L◌│◌│◌◌║ 7} \\
\texttt{╚══╧═╧══╝} \\
\\ 
\texttt{```} \\
            }
        }
    }
    & & \\ \\

    \theutterance \stepcounter{utterance}  
    & & & \multicolumn{2}{p{0.3\linewidth}}{
        \cellcolor[rgb]{0.9,0.9,0.9}{
            \makecell[{{p{\linewidth}}}]{
                \texttt{\tiny{[GM$|$GM]}}
                \texttt{End states:} \\
\texttt{```} \\
\texttt{Player 1:} \\
\texttt{ 1234567} \\
\texttt{╔════╤══╗} \\
\texttt{║◌◌◌◌│◌◌║ 1} \\
\texttt{║◌◌◌◌├──╢ 2} \\
\texttt{║◌◌CP│◌◌║ 3} \\
\texttt{║◌┌──┼──╢ 4} \\
\texttt{║◌│◌L│◌◌║ 5} \\
\texttt{╟─┴┬─┤◌◌║ 6} \\
\texttt{║◌◌│◌│◌◌║ 7} \\
\texttt{╚══╧═╧══╝} \\
\\ 
\\ 
\texttt{Player 2:} \\
\texttt{ 1234567} \\
\texttt{╔════╤══╗} \\
\texttt{║◌◌◌◌│◌◌║ 1} \\
\texttt{║◌◌◌◌├──╢ 2} \\
\texttt{║◌◌CP│◌◌║ 3} \\
\texttt{║◌┌──┼──╢ 4} \\
\texttt{║◌│◌L│◌◌║ 5} \\
\texttt{╟─┴┬─┤◌◌║ 6} \\
\texttt{║◌◌│◌│◌◌║ 7} \\
\texttt{╚══╧═╧══╝} \\
\\ 
\texttt{```} \\
            }
        }
    }
    & & \\ \\

    \theutterance \stepcounter{utterance}  
    & & & \multicolumn{2}{p{0.3\linewidth}}{
        \cellcolor[rgb]{0.9,0.9,0.9}{
            \makecell[{{p{\linewidth}}}]{
                \texttt{\tiny{[GM$|$GM]}}
                \texttt{* success: False} \\
\texttt{* lose: False} \\
\texttt{* aborted: True} \\
\texttt{{-}{-}{-}{-}{-}{-}{-}} \\
\texttt{* Shifts: 2.00} \\
\texttt{* Max Shifts: 4.00} \\
\texttt{* Min Shifts: 2.00} \\
\texttt{* End Distance Sum: 0.00} \\
\texttt{* Init Distance Sum: 9.04} \\
\texttt{* Expected Distance Sum: 12.57} \\
\texttt{* Penalties: 9.00} \\
\texttt{* Max Penalties: 8.00} \\
\texttt{* Rounds: 8.00} \\
\texttt{* Max Rounds: 12.00} \\
\texttt{* Object Count: 3.00} \\
            }
        }
    }
    & & \\ \\

\end{supertabular}
}

\end{document}
